\section*{ID10073: Gauge Symmetry: G(t)}
\paragraph{Metadata.}
PiID: 0435220277098 | RTEID: RTE:P39115.6960:T1.6188 | UUID: b6501a69-5c28-55ff-a0fe-1878c13bd70d

\paragraph{Canonical Statement.}
\textbackslash{}begin\{definition\}[Gauge Transformation] A \textbackslash{}emph\{gauge transformation\} is a natural transformation from the identity functor on state space to itself that acts fiberwise by the gauge group \textbackslash{}(U(1)\textbackslash{}times\textbackslash{}mathrm\{SO\}(3)\textbackslash{}). For a state \textbackslash{}(\textbackslash{}Psi(t)\textbackslash{}), the component at time \textbackslash{}(t\textbackslash{}) is given by \textbackslash{}[ \textbackslash{}tau\_t: \textbackslash{}Psi(t) \textbackslash{}longmapsto g(t)\textbackslash{}cdot\textbackslash{}Psi(t), \textbackslash{}] where \textbackslash{}(g(t) = (\textbackslash{}mathrm\{e\}\textasciicircum{}\{\textbackslash{}mathrm\{i\}\textbackslash{}theta(t)\}, R(\textbackslash{}alpha(t),\textbackslash{}beta(t),\textbackslash{}gamma(t)))\textbackslash{}). \textbackslash{}end\{definition\}

\paragraph{Canonical Equation.}
\[
g(t) = (\mathrm{e}^{\mathrm{i}\theta(t)}, R(\alpha(t),\beta(t),\gamma(t)))
\]

\paragraph{Dictionary.}
Gauge Symmetry: G(t) — g(t) = (e\textasciicircum{}iθ(t), R(α(t),β(t),γ(t)))

\paragraph{Encyclopedia.}
Gauge Symmetry: G(t) is an algebraic definition, written g(t) = (e\textasciicircum{}iθ(t), R(α(t),β(t),γ(t))). It is an algebraic definition expressing the quantity directly in terms of the variables on its right-hand side. It is built from θ, α, β, γ, R, defining g(t). Within the Trawin Zero Point registry it serves as one element of the framework's mathematical narrative.
